
\subsubsection{Torque Analysis}

We selected high-torque servos, so that body weight does not become a binding design constraint.
To quantify the actual limits this gives us, we invert the static torque equations and solve for the \textbf{maximum body weight} the leg can support as a function of link geometry and joint angles.
We model a single leg as a 3-link serial chain rotating in the sagittal (pitch) plane, like a finger.

% TODO
% Link ./img/robot-leg-figure.png and ./img/robot-leg-figure-bent.png
% showign side by side how we simulated it, to show how we formulate the problem.

\vspace{0.5cm}
\textbf{Leg model.}
Each leg consists of three links: $L_1$ (shoulder to hip joint, length $r_1$), $L_2$ (hip to knee, length $r_2$), and $L_3$ (knee to foot tip, length $r_3$).
The hip and knee motors act on pitch axes and bear static gravity torque, while the shoulder motor acts on a yaw axis (parallel to gravity) and bears only dynamic rotational inertia during swing phase.

\vspace{0.5cm}
\textbf{Knee motor torque.}
The knee servo supports link $L_3$ and the ground-reaction tip load $m_{\text{tip}}$.
All relevant masses lie along $L_3$, so a single angle $\theta_{\text{knee}}$ governs the gravity projection:
%
\begin{align*}
  \tau_{L_3}
  &= \left( m_{L_3} \cdot \frac{r_3}{2}
  + m_{\text{tip}} \cdot r_3 \right)
  \cos\theta_{\text{knee}}
\end{align*}

\textbf{Hip motor torque.}
The hip servo supports $L_2$, the knee servo mass, $L_3$, and the tip load.
Because the hip and knee angles are independent, masses on $L_2$ project by
$\cos\theta_{\text{hip}}$ while masses on $L_3$ project by
$\cos\theta_{\text{knee}}$:
%
\begin{align*}
  \tau_{L_2}
  &= m_{L_2}\,\frac{r_2}{2}\,\cos\theta_{\text{hip}}
  + m_{\text{servo,knee}}\,r_2\,\cos\theta_{\text{hip}} \\
  &\quad + m_{L_3}\!\left(r_2\cos\theta_{\text{hip}}
  + \frac{r_3}{2}\cos\theta_{\text{knee}}\right) \\
  &\quad + m_{\text{tip}}\!\left(r_2\cos\theta_{\text{hip}}
  + r_3\cos\theta_{\text{knee}}\right)
\end{align*}

\textbf{Inverse solver - weight budget.}
With the motor choice settled, the practical question is how much total mass we can allocate to the chassis, electronics, and payload.
Setting each torque equal to the rated stall torque and solving for $m_{\text{tip}}$ yields the per-leg load capacity.
For the hip:
%
\begin{align*}
  m_{\text{tip,max}}^{\text{hip}}
  &= \frac{
    \tau_{\text{rated}}
    - \bigl(m_{L_2}\,\tfrac{r_2}{2}
    + m_{\text{servo,knee}}\,r_2\bigr)\cos\theta_{\text{hip}}
    - m_{L_3}\bigl(r_2\cos\theta_{\text{hip}}
    + \tfrac{r_3}{2}\cos\theta_{\text{knee}}\bigr)
  }{
    r_2\cos\theta_{\text{hip}} + r_3\cos\theta_{\text{knee}}
  }
\end{align*}
%
The binding constraint is whichever servo saturates first.  
Total supportable body mass is then:
%
\[
  m_{\text{body}}
  = \min\!\bigl(m_{\text{tip,max}}^{\text{knee}},\;
  m_{\text{tip,max}}^{\text{hip}}\bigr)
  \times n_{\text{legs in stance}}
\]

\vspace{0.5cm}
\textbf{Dynamic yaw torque.}
As a secondary analysis, we compute the moment of inertia of the fully extended leg about the shoulder yaw axis (using the parallel-axis theorem for each link and point mass) and derive the required yaw torque $\tau_{\text{yaw}} = I_{\text{leg}} \cdot \alpha$ for a target swing-phase angular acceleration $\alpha$.
% TODO: report numerical I_leg and required yaw torque

\vspace{0.5cm}
An interactive Jupyter notebook (\texttt{torque\_calculations.ipynb}) allows exploration of these equations across the full range of joint angles and design parameters.

% Link the image ./img/interactive-figure-image.png showing the actual interactive simulator.
